\section{Performance Engineering}
Charge deposition is memory-bound on CPU because particles are read sequentially while grid writes are irregular and often non-local.
We apply three complementary optimizations: cell sorting for spatial locality, SoA layout for vector-friendly particle reads, and OpenMP grid privatization to avoid atomic contention.

\subsection{Particle Sorting}
Before deposition, particles are sorted by cell index $c = i_y N_x + i_x$, computed from periodic-wrapped coordinates.
For AoS, \texttt{std::sort} reorders the particle array in place; for SoA, an index permutation is sorted and attributes are gathered into a new buffer.
Sorting costs $\mathcal{O}(N_p \log N_p)$ but clusters particles that write to nearby grid cells, improving cache reuse during the scatter loop.
In timed microbenchmarks, sort and deposit are measured separately so that locality benefits are not conflated with reordering overhead.

\subsection{AoS vs SoA}
AoS stores $(x,y,v_x,v_y,q,m)$ in a 48-byte struct per particle; SoA stores each attribute in a separate contiguous array (40 bytes read per particle during deposition).
SoA enables unit-stride access to positions and charges, which is favorable when the inner loop touches only a subset of attributes.
AoS can be faster to sort because reordering a single array avoids the SoA gather pass.
Our benchmarks compare both layouts under identical schemes and thread counts.

\subsection{OpenMP Privatization}
Parallel CPU deposition partitions the particle index range across $T$ threads.
Each thread allocates a private copy of the full $\rho$ grid, deposits its particle subset locally, then merges into the global grid inside an \texttt{omp critical} region:
\begin{equation}
    \rho^{\mathrm{global}}_k \mathrel{+}= \sum_{t=0}^{T-1} \rho^{\mathrm{local},t}_k.
\end{equation}
This avoids per-cell atomics that serialize cache lines under high thread counts.
The trade-off is $T\times$ temporary grid storage and a merge pass whose cost grows with grid size and thread count.

\subsection{GPU Strategies}
CUDA kernels implementing atomic and privatized-tile deposition compile when \texttt{BUILD\_CUDA=ON}.
The atomic variant launches one thread per particle and scatters into the global grid with hardware \texttt{atomicAdd}.
The privatized variant launches one block per $16\times16$ spatial tile; threads stride over all particles, accumulate matching stencil contributions in shared memory with \texttt{atomicAdd}, then merge each tile into global memory.
This design preserves charge conservation but costs $\mathcal{O}(N_{\mathrm{tiles}}\,N_p)$ work because every tile block scans the full particle list.
Atomic GPU deposition is therefore the preferred production backend on Tesla T4; privatized tiles are retained for correctness validation and as a baseline for future particle--tile binning.
Both backends share the same \texttt{DepositionConfig} interface as the CPU path.
